\documentclass{article}


\usepackage{arxiv}

\usepackage[utf8]{inputenc} % allow utf-8 input
\usepackage[spanish]{babel}
\usepackage[T1]{fontenc}    % use 8-bit T1 fonts
\usepackage{hyperref}       % hyperlinks
\usepackage{url}            % simple URL typesetting
\usepackage{booktabs}       % professional-quality tables
\usepackage{amsfonts}       % blackboard math symbols
\usepackage{nicefrac}       % compact symbols for 1/2, etc.
\usepackage{microtype}      % microtypography
\usepackage{lipsum}
\usepackage{graphicx}
\usepackage{fancyhdr}
\usepackage{float}


\title{TP 1.1 Simulación de una ruleta}


\author{
 Tomás Lardizábal \\
  Legajo 47433\\
  Universidad Tecnológica Nacional - FRRO\\
  Zeballos 1341, S2000, Argentina \\
  \texttt{tlardizabal1@gmail.com} \\
  %% examples of more authors
   \And
 Iñaki Díaz\\
  Legajo 48944\\
  Universidad Tecnológica Nacional - FRRO\\
  Zeballos 1341, S2000, Argentina \\
  \texttt{inakidiaz2808@gmail.com} \\
    \And
 Tomás Splivalo\\
  Legajo 51665 \\
  Universidad Tecnológica Nacional - FRRO\\
  Zeballos 1341, S2000, Argentina \\
  \texttt{tomassplivalo1@gmail.com} \\
    \And
 Luciano Armas\\
  Legajo 47181\\
  Universidad Tecnológica Nacional - FRRO\\
  Zeballos 1341, S2000, Argentina \\
  \texttt{lucianoarmas11@gmail.com} \\
}

\date{\today}
\setlength{\itemsep}{0pt}
\begin{document}
\maketitle

\pagestyle{fancy}
\fancyhead{} % Elimina encabezado predeterminado
\rhead{Simulación de una ruleta, \today}

\begin{abstract}
Mediante el siguiente estudio se pretende realizar un análisis y evaluación de las distintas estrategias de juego en la ruleta a través de simulaciones en lenguaje Python, contemplando dos escenarios distintos: 
\begin{itemize}
  \item Uno en el cual el jugador posee un capital infinito
  \item Otro escenario en el que el jugador posee un capital acotado.
\end{itemize}
 El objetivo es conocer si hay posibilidades de lograr resultados positivos que hagan crecer el capital apostado.
Se evalúan las estrategias Martingala, Martingala Invertida, Fibonacci y D'Alembert.
\end{abstract}


\section{Introducción}
El juego de la ruleta ha sido analizado desde múltiples enfoques, tanto lúdicos como científicos, para responder a la pregunta de si existe algún algoritmo que pueda predecir la aleatoriedad de este juego, permitiéndonos obtener beneficios.
En este informe pondremos a prueba algunas de las famosas estrategias de apuestas para analizar su comportamiento a lo largo de múltiples tiradas y estudiar las probabilidades de obtener ganancia con cada estrategia, tomando una ruleta europea ideal.




\section{Metodología}
El estudio lleva a cabo una simulación del juego de la ruleta desarrollada en el lenguaje Python. En dicha simulación se aplican las siguientes técnicas de apuesta:
\begin{itemize}
  \item Apuesta a un único numero. Dicha técnica tiene un beneficio de 35 a 1, es decir, en caso de acertar el resultado al jugador se le paga lo que apostó multiplicado por 35.
  \item Todo negro. Dicha técnica tiene un beneficio 1 a 1, es decir, en caso de acertar el resultado al jugador se le paga lo que apostó.
\end{itemize} 
Para ambas técnicas se utilizan distintas estrategias conocidas en el ambiente del juego.
Cada estrategia será puesta en dos escenarios distintos, uno en donde el apostador posee un capital ilimitado y otro en el que el mismo es acotado. Posteriormente se analizaran los resultados en cada escenario.





\section{Estrategias}
\subsection{Martingala}
El método Martingala es un sistema de apuestas que se basa en duplicar la apuesta después de cada pérdida con la esperanza de recuperar todas las pérdidas con una sola victoria. 
Las desventajas de este metodo son: la necesidad de un capital considerable, la existencia de límites de mesa en los casinos y la ventaja inherente de la casa. Además, tenemos que tener en cuenta “la falacia del jugador”, una creencia de que las pérdidas pasadas aumentan la probabilidad de una futura victoria, que en metodologías como la martingala puede ser una premisa que te lleve a la bancarrota.
A pesar de que es posible obtener ganancias a corto plazo utilizando el sistema Martingala, no es una estrategia confiable ni sostenible para ganar en la ruleta a largo plazo. El riesgo de sufrir pérdidas sustanciales debido a rachas de pérdidas inevitables supera con creces la pequeña ganancia que se obtiene con cada secuencia exitosa. 
En última instancia, el sistema Martingala, a pesar de su popularidad y su aparente lógica, presenta fallas fundamentales debido a las restricciones prácticas del juego en casinos y a la propia probabilidad de la ruleta, que mostraremos en este trabajo.


\subsection{Martingala Inversa o Antimartingala}
La Martingala Inversa (o Antimartingala) es un sistema de apuestas que crea un patrón en el que el tamaño de la apuesta está directamente relacionado con el resultado del juego.
Según este sistema, un jugador debería:
\begin{itemize}
  \item Duplicar la apuesta después de una victoria.
  \item Volver a la apuesta inicial después de una pérdida.
\end{itemize}
Este sistema se basa en la premisa de que los jugadores de ruleta pueden sacar provecho de las rachas ganadoras con la estrategia correcta de gestión del bankroll (monto de dinero que un jugador está dispuesto a arriesgar), y que pueden reducir el riesgo de perder su dinero cuando las apuestas no van a su favor.
Esta estrategia, que es una variación del sistema Martingala mencionado en el punto anterior, surgió de la creencia de que tanto las pérdidas como las ganancias tienden a producirse en serie.

Con la Antimartingala, los jugadores de ruleta intentan maximizar una parte de las ganancias durante las rachas positivas, reduciéndolas inmediatamente durante las rachas negativas para intentar preservar su equilibrio.

\subsection{Fibonacci}
El sistema Fibonacci es una estrategia de apuestas progresiva que se basa en la famosa secuencia matemática de Fibonacci: [1,1,2,3,5,8,13,21,…,] en la que cada número es la suma de los dos anteriores. Este sistema fue adaptado al juego con el objetivo de recuperar pérdidas de forma más moderada que con estrategias más agresivas como la Martingala. La mecánica del sistema es sencilla: el jugador comienza apostando una unidad (por ejemplo, \$1).

\begin{itemize}
    \item Si pierde, avanza un paso en la secuencia y realiza la siguiente apuesta con ese nuevo valor.
    \item Si gana, retrocede dos pasos en la secuencia y apuesta el valor correspondiente.
    \item Si está en uno de los primeros dos números de la secuencia (los dos “1”), simplemente permanece en ese punto o vuelve a iniciar.
\end{itemize}

Esta estrategia tiene como objetivo recuperar gradualmente las pérdidas acumuladas con una estructura de progresión más conservadora que reduce el riesgo de grandes pérdidas en comparación con Martingala. Sin embargo, sigue requiriendo un capital suficientemente amplio para soportar una posible racha de pérdidas. La ventaja principal del método es que limita el ritmo de aumento de las apuestas, haciendo menos probable que el jugador se vea rápidamente afectado por los límites de la mesa o por quedarse sin fondos. No obstante, como cualquier estrategia basada en progresiones, no altera la ventaja matemática de la casa ni garantiza beneficios a largo plazo. En resumen, aunque el sistema Fibonacci puede ofrecer una experiencia de juego más estable y manejable que estrategias más agresivas, sigue estando sujeto a las mismas limitaciones inherentes a todos los sistemas de apuestas: la independencia de los eventos en la ruleta, el margen de la casa y la posibilidad de enfrentar secuencias de pérdidas prolongadas.

\subsection{D'Alembert}
El sistema D’Alembert es una estrategia de apuestas utilizada comúnmente en juegos de azar como la ruleta, y se basa en un principio de progresión aritmética. A diferencia de sistemas más agresivos como la Martingala, la estrategia D’Alembert propone una progresión mucho más moderada, lo que la convierte en una de las metodologías más simples y conservadoras dentro del análisis de apuestas.
Esta estrategia está pensada para aplicarse en apuestas de tipo binario, también conocidas como apuestas "simples", que tienen una probabilidad cercana al 50\%. Entre ellas se incluyen Rojo/Negro, Par/Impar, y Bajo/Alto (1–18 / 19–36). La lógica detrás de este sistema parte del supuesto de que, a largo plazo, los resultados tienden a equilibrarse, alternándose las rachas positivas y negativas.
El funcionamiento del sistema es el siguiente: el apostador define una "unidad base" que representa el monto mínimo de apuesta. Esta unidad debe seleccionarse de manera responsable, idealmente representando entre un 0,33\%y un 1\% del capital total disponible, para evitar riesgos excesivos en caso de una racha prolongada de pérdidas.
La dinámica consiste en aumentar la apuesta en una unidad tras cada pérdida y disminuirla en una unidad tras cada ganancia. Por ejemplo, si se comienza apostando una unidad y se pierde, la siguiente apuesta será de dos unidades; si se gana, se reduce a una. Este patrón se mantiene durante toda la sesión de juego.
El objetivo del sistema es aprovechar los momentos en que la secuencia de pérdidas es compensada por una sucesión posterior de ganancias, permitiendo recuperar paulatinamente las pérdidas acumuladas. No obstante, como toda estrategia de juego basada en probabilidades independientes, el sistema D’Alembert no garantiza beneficios a largo plazo, aunque sí ofrece una mayor estabilidad y menor volatilidad que otras estrategias más agresivas.




\section{Casos de estudio}
Tomaremos una cantidad máxima de 10.000 tiradas de la ruleta, ya que consideramos que es más que suficiente para realizar el análisis. En el caso de capital finito, la simulación terminará cuando se llegue a la bancarrota.
Cabe aclarar que en todos los giros de ruleta en los que se utilice la técnica de "numero único" se elegirá el numero 11. Los números de la técnica "todos negros" estan compuesto por: 2, 4, 6, 8, 10, 11, 13, 15, 17, 20, 22, 24, 26, 28, 29, 31, 33, 35.


\subsection{Capital finito}
Para la simulación con capital finito se utilizará un capital de 10.000 unidades con apuesta inicial de 10 que se irá modificando por las distintas estrategias. La simulación finalizará al llegar a la quiebra o a las 10.000 tiradas a la ruleta.


\subsubsection{D'Alembert - Técnica - Único numero}
 \begin{figure}[H]
   \centering
    \includegraphics[width=0.8\linewidth]{graficos/capital_finito/dalembert_numero_unico/flujo_caja_d_f.png}
    \caption{Flujo de caja - 1° corrida}
    \label{fig:flujo_caja_d_f-4}
\end{figure}
 \begin{figure}[H]
   \centering
    \includegraphics[width=0.8\linewidth]{graficos/capital_finito/dalembert_numero_unico/corridas_5_d_f.png}
    \caption{Flujo de caja - 5 corridas}
    \label{fig:corridas_5_d_f-4}
\end{figure}
 \begin{figure}[H]
   \centering
    \includegraphics[width=0.8\linewidth]{graficos/capital_finito/dalembert_numero_unico/frecuencia_relativa_evolucion_d_f.png}
    \caption{Histograma de frecuencia}
    \label{fig:frecuencia_relativa_evolucion_d_f-4}
\end{figure}
 \begin{figure}[H]
   \centering
    \includegraphics[width=0.8\linewidth]{graficos/capital_finito/dalembert_numero_unico/frecuencias_multiples_d_f.png}
    \caption{Histograma de frecuencia}
    \label{fig:frecuencias_multiples_d_f-4}
\end{figure}


\subsubsection{Fibonacci - Técnica - Único numero}
 \begin{figure}[H]
   \centering
    \includegraphics[width=0.8\linewidth]{graficos/capital_finito/fibonacci_numero_unico/flujo_caja_f_f.png}
    \caption{Flujo de caja - 1° corrida}
    \label{fig:flujo_caja_f_f-4}
\end{figure}
 \begin{figure}[H]
   \centering
    \includegraphics[width=0.8\linewidth]{graficos/capital_finito/fibonacci_numero_unico/corridas_5_f_f.png}
    \caption{Flujo de caja - 5 corridas}
    \label{fig:corridas_5_f_f-4}
\end{figure}
 \begin{figure}[H]
   \centering
    \includegraphics[width=0.8\linewidth]{graficos/capital_finito/fibonacci_numero_unico/frecuencia_relativa_evolucion_f_f.png}
    \caption{Histograma de frecuencia}
    \label{fig:frecuencia_relativa_evolucion_f_f-4}
\end{figure}
 \begin{figure}[H]
   \centering
    \includegraphics[width=0.8\linewidth]{graficos/capital_finito/fibonacci_numero_unico/frecuencias_multiples_f_f.png}
    \caption{Histograma de frecuencia}
    \label{fig:frecuencias_multiples_f_f-4}
\end{figure}



\subsubsection{Martingala - Técnica - Único numero}
 \begin{figure}[H]
   \centering
    \includegraphics[width=0.8\linewidth]{graficos/capital_finito/martingala_numero_unico/flujo_caja_m_f.png}
    \caption{Flujo de caja - 1° corrida}
    \label{fig:flujo_caja_m_f-4}
\end{figure}
 \begin{figure}[H]
   \centering
    \includegraphics[width=0.8\linewidth]{graficos/capital_finito/martingala_numero_unico/corridas_5_m_f.png}
    \caption{Flujo de caja - 5 corridas}
    \label{fig:corridas_5_m_f-4}
\end{figure}
 \begin{figure}[H]
   \centering
    \includegraphics[width=0.8\linewidth]{graficos/capital_finito/martingala_numero_unico/frecuencia_relativa_evolucion_m_f.png}
    \caption{Histograma de frecuencia}
    \label{fig:frecuencia_relativa_evolucion_m_f-4}
\end{figure}
 \begin{figure}[H]
   \centering
    \includegraphics[width=0.8\linewidth]{graficos/capital_finito/martingala_numero_unico/frecuencias_multiples_m_f.png}
    \caption{Histograma de frecuencia}
    \label{fig:frecuencias_multiples_m_f-4}
\end{figure}



\subsubsection{Martingala Invertido - Técnica - Único numero}
 \begin{figure}[H]
   \centering
    \includegraphics[width=0.8\linewidth]{graficos/capital_finito/martingala_invertida_numero_unico/flujo_caja_o_f.png}
    \caption{Flujo de caja - 1° corrida}
    \label{fig:flujo_caja_mi_f-4}
\end{figure}
 \begin{figure}[H]
   \centering
    \includegraphics[width=0.8\linewidth]{graficos/capital_finito/martingala_invertida_numero_unico/corridas_5_o_f.png}
    \caption{Flujo de caja - 5 corridas}
    \label{fig:corridas_5_mi_f-4}
\end{figure}
 \begin{figure}[H]
   \centering
    \includegraphics[width=0.8\linewidth]{graficos/capital_finito/martingala_invertida_numero_unico/frecuencia_relativa_evolucion_o_f.png}
    \caption{Histograma de frecuencia}
    \label{fig:frecuencia_relativa_evolucion_mi_f-4}
\end{figure}
 \begin{figure}[H]
   \centering
    \includegraphics[width=0.8\linewidth]{graficos/capital_finito/martingala_invertida_numero_unico/frecuencias_multiples_o_f.png}
    \caption{Histograma de frecuencia}
    \label{fig:frecuencias_multiples_mi_f-4}
\end{figure}






\subsubsection{D'Alembert - Técnica - Números negros}
 \begin{figure}[H]
   \centering
    \includegraphics[width=0.8\linewidth]{graficos/capital_finito/dalembert_numeros_negros/flujo_caja_d_f.png}
    \caption{Flujo de caja - 1° corrida}
    \label{fig:flujo_caja_d_f_numnegro-4}
\end{figure}
 \begin{figure}[H]
   \centering
    \includegraphics[width=0.8\linewidth]{graficos/capital_finito/dalembert_numeros_negros/corridas_5_d_f.png}
    \caption{Flujo de caja - 5 corridas}
    \label{fig:corridas_5_d_f_numnegro-4}
\end{figure}
 \begin{figure}[H]
   \centering
    \includegraphics[width=0.8\linewidth]{graficos/capital_finito/dalembert_numeros_negros/frecuencia_relativa_evolucion_d_f.png}
    \caption{Histograma de frecuencia}
    \label{fig:frecuencia_relativa_evolucion_d_f_numnegro-4}
\end{figure}
 \begin{figure}[H]
   \centering
    \includegraphics[width=0.8\linewidth]{graficos/capital_finito/dalembert_numeros_negros/frecuencias_multiples_d_f.png}
    \caption{Histograma de frecuencia}
    \label{fig:frecuencias_multiples_d_f_numnegro-4}
\end{figure}


\subsubsection{Fibonacci - Técnica - Números negros}
 \begin{figure}[H]
   \centering
    \includegraphics[width=0.8\linewidth]{graficos/capital_finito/fibonacci_numeros_negros/flujo_caja_f_f.png}
    \caption{Flujo de caja - 1° corrida}
    \label{fig:flujo_caja_f_f_numnegro-4}
\end{figure}
 \begin{figure}[H]
   \centering
    \includegraphics[width=0.8\linewidth]{graficos/capital_finito/fibonacci_numeros_negros/corridas_5_f_f.png}
    \caption{Flujo de caja - 5 corridas}
    \label{fig:corridas_5_f_f_numnegro-4}
\end{figure}
 \begin{figure}[H]
   \centering
    \includegraphics[width=0.8\linewidth]{graficos/capital_finito/fibonacci_numeros_negros/frecuencia_relativa_evolucion_f_f.png}
    \caption{Histograma de frecuencia}
    \label{fig:frecuencia_relativa_evolucion_f_f_numnegro-4}
\end{figure}
 \begin{figure}[H]
   \centering
    \includegraphics[width=0.8\linewidth]{graficos/capital_finito/fibonacci_numeros_negros/frecuencias_multiples_f_f.png}
    \caption{Histograma de frecuencia}
    \label{fig:frecuencias_multiples_f_f_numnegro-4}
\end{figure}



\subsubsection{Martingala - Técnica - Números negros}
 \begin{figure}[H]
   \centering
    \includegraphics[width=0.8\linewidth]{graficos/capital_finito/martingala_numeros_negros/flujo_caja_m_f.png}
    \caption{Flujo de caja - 1° corrida}
    \label{fig:flujo_caja_m_f_numnegro-4}
\end{figure}
 \begin{figure}[H]
   \centering
    \includegraphics[width=0.8\linewidth]{graficos/capital_finito/martingala_numeros_negros/corridas_5_m_f.png}
    \caption{Flujo de caja - 5 corridas}
    \label{fig:corridas_5_m_f_numnegro-4}
\end{figure}
 \begin{figure}[H]
   \centering
    \includegraphics[width=0.8\linewidth]{graficos/capital_finito/martingala_numeros_negros/frecuencia_relativa_evolucion_m_f.png}
    \caption{Histograma de frecuencia}
    \label{fig:frecuencia_relativa_evolucion_m_f_numnegro-4}
\end{figure}
 \begin{figure}[H]
   \centering
    \includegraphics[width=0.8\linewidth]{graficos/capital_finito/martingala_numeros_negros/frecuencias_multiples_m_f.png}
    \caption{Histograma de frecuencia}
    \label{fig:frecuencias_multiples_m_f_numnegro-4}
\end{figure}



\subsubsection{Martingala Invertido - Técnica - Números negros}
 \begin{figure}[H]
   \centering
    \includegraphics[width=0.8\linewidth]{graficos/capital_finito/martingala_invertida_numeros_negros/flujo_caja_o_f.png}
    \caption{Flujo de caja - 1° corrida}
    \label{fig:flujo_caja_mi_f_numnegro-4}
\end{figure}
 \begin{figure}[H]
   \centering
    \includegraphics[width=0.8\linewidth]{graficos/capital_finito/martingala_invertida_numeros_negros/corridas_5_o_f.png}
    \caption{Flujo de caja - 5 corridas}
    \label{fig:corridas_5_mi_f_numnegro-4}
\end{figure}
 \begin{figure}[H]
   \centering
    \includegraphics[width=0.8\linewidth]{graficos/capital_finito/martingala_invertida_numeros_negros/frecuencia_relativa_evolucion_o_f.png}
    \caption{Histograma de frecuencia}
    \label{fig:frecuencia_relativa_evolucion_mi_f_numnegro-4}
\end{figure}
 \begin{figure}[H]
   \centering
    \includegraphics[width=0.8\linewidth]{graficos/capital_finito/martingala_invertida_numeros_negros/frecuencias_multiples_o_f.png}
    \caption{Histograma de frecuencia}
    \label{fig:frecuencias_multiples_mi_f_numnegro-4}
\end{figure}






\subsection{Capital infinito}
Para la simulación con capital infinito se realziará una apuesta inicial de 10 que se irá modificando por las distintas estrategias. La simulación finalizará al llegar a las 10.000 tiradas a la ruleta.


\subsubsection{D'Alembert - Técnica - Único numero}
 \begin{figure}[H]
   \centering
    \includegraphics[width=0.8\linewidth]{graficos/capital_infinito/dalembert_numero_unico/flujo_caja_d_i.png}
    \caption{Flujo de caja - 1° corrida}
    \label{fig:flujo_caja_d_i-4}
\end{figure}
 \begin{figure}[H]
   \centering
    \includegraphics[width=0.8\linewidth]{graficos/capital_infinito/dalembert_numero_unico/corridas_5_d_i.png}
    \caption{Flujo de caja - 5 corridas}
    \label{fig:corridas_5_d_i-4}
\end{figure}
 \begin{figure}[H]
   \centering
    \includegraphics[width=0.8\linewidth]{graficos/capital_infinito/dalembert_numero_unico/frecuencia_relativa_evolucion_d_i.png}
    \caption{Histograma de frecuencia}
    \label{fig:frecuencia_relativa_evolucion_d_i-4}
\end{figure}
 \begin{figure}[H]
   \centering
    \includegraphics[width=0.8\linewidth]{graficos/capital_infinito/dalembert_numero_unico/frecuencias_multiples_d_i.png}
    \caption{Histograma de frecuencia}
    \label{fig:frecuencias_multiples_d_i-4}
\end{figure}


\subsubsection{Fibonacci - Técnica - Único numero}
 \begin{figure}[H]
   \centering
    \includegraphics[width=0.8\linewidth]{graficos/capital_infinito/fibonacci_numero_unico/flujo_caja_f_i.png}
    \caption{Flujo de caja - 1° corrida}
    \label{fig:flujo_caja_f_i-4}
\end{figure}
 \begin{figure}[H]
   \centering
    \includegraphics[width=0.8\linewidth]{graficos/capital_infinito/fibonacci_numero_unico/corridas_5_f_i.png}
    \caption{Flujo de caja - 5 corridas}
    \label{fig:corridas_5_f_i-4}
\end{figure}
 \begin{figure}[H]
   \centering
    \includegraphics[width=0.8\linewidth]{graficos/capital_infinito/fibonacci_numero_unico/frecuencia_relativa_evolucion_f_i.png}
    \caption{Histograma de frecuencia}
    \label{fig:frecuencia_relativa_evolucion_f_i-4}
\end{figure}
 \begin{figure}[H]
   \centering
    \includegraphics[width=0.8\linewidth]{graficos/capital_infinito/fibonacci_numero_unico/frecuencias_multiples_f_i.png}
    \caption{Histograma de frecuencia}
    \label{fig:frecuencias_multiples_f_i-4}
\end{figure}



\subsubsection{Martingala - Técnica - Único numero}
 \begin{figure}[H]
   \centering
    \includegraphics[width=0.8\linewidth]{graficos/capital_infinito/martingala_numero_unico/flujo_caja_m_i.png}
    \caption{Flujo de caja - 1° corrida}
    \label{fig:flujo_caja_m_i-4}
\end{figure}
 \begin{figure}[H]
   \centering
    \includegraphics[width=0.8\linewidth]{graficos/capital_infinito/martingala_numero_unico/corridas_5_m_i.png}
    \caption{Flujo de caja - 5 corridas}
    \label{fig:corridas_5_m_i-4}
\end{figure}
 \begin{figure}[H]
   \centering
    \includegraphics[width=0.8\linewidth]{graficos/capital_infinito/martingala_numero_unico/frecuencia_relativa_evolucion_m_i.png}
    \caption{Histograma de frecuencia}
    \label{fig:frecuencia_relativa_evolucion_m_i-4}
\end{figure}
 \begin{figure}[H]
   \centering
    \includegraphics[width=0.8\linewidth]{graficos/capital_infinito/martingala_numero_unico/frecuencias_multiples_m_i.png}
    \caption{Histograma de frecuencia}
    \label{fig:frecuencias_multiples_m_i-4}
\end{figure}



\subsubsection{Martingala Invertido - Técnica - Único numero}
 \begin{figure}[H]
   \centering
    \includegraphics[width=0.8\linewidth]{graficos/capital_infinito/martingala_invertida_numero_unico/flujo_caja_o_i.png}
    \caption{Flujo de caja - 1° corrida}
    \label{fig:flujo_caja_mi_i-4}
\end{figure}
 \begin{figure}[H]
   \centering
    \includegraphics[width=0.8\linewidth]{graficos/capital_infinito/martingala_invertida_numero_unico/corridas_5_o_i.png}
    \caption{Flujo de caja - 5 corridas}
    \label{fig:corridas_5_mi_i-4}
\end{figure}
 \begin{figure}[H]
   \centering
    \includegraphics[width=0.8\linewidth]{graficos/capital_infinito/martingala_invertida_numero_unico/frecuencia_relativa_evolucion_o_i.png}
    \caption{Histograma de frecuencia}
    \label{fig:frecuencia_relativa_evolucion_mi_i-4}
\end{figure}
 \begin{figure}[H]
   \centering
    \includegraphics[width=0.8\linewidth]{graficos/capital_infinito/martingala_invertida_numero_unico/frecuencias_multiples_o_i.png}
    \caption{Histograma de frecuencia}
    \label{fig:frecuencias_multiples_mi_i-4}
\end{figure}






\subsubsection{D'Alembert - Técnica - Números negros}
 \begin{figure}[H]
   \centering
    \includegraphics[width=0.8\linewidth]{graficos/capital_infinito/dalembert_numeros_negros/flujo_caja_d_i.png}
    \caption{Flujo de caja - 1° corrida}
    \label{fig:flujo_caja_d_i_numnegro-4}
\end{figure}
 \begin{figure}[H]
   \centering
    \includegraphics[width=0.8\linewidth]{graficos/capital_infinito/dalembert_numeros_negros/corridas_5_d_i.png}
    \caption{Flujo de caja - 5 corridas}
    \label{fig:corridas_5_d_i_numnegro-4}
\end{figure}
 \begin{figure}[H]
   \centering
    \includegraphics[width=0.8\linewidth]{graficos/capital_infinito/dalembert_numeros_negros/frecuencia_relativa_evolucion_d_i.png}
    \caption{Histograma de frecuencia}
    \label{fig:frecuencia_relativa_evolucion_d_i_numnegro-4}
\end{figure}
 \begin{figure}[H]
   \centering
    \includegraphics[width=0.8\linewidth]{graficos/capital_infinito/dalembert_numeros_negros/frecuencias_multiples_d_i.png}
    \caption{Histograma de frecuencia}
    \label{fig:frecuencias_multiples_d_i_numnegro-4}
\end{figure}


\subsubsection{Fibonacci - Técnica - Números negros}
 \begin{figure}[H]
   \centering
    \includegraphics[width=0.8\linewidth]{graficos/capital_infinito/fibonacci_numeros_negros/flujo_caja_f_i.png}
    \caption{Flujo de caja - 1° corrida}
    \label{fig:flujo_caja_f_i_numnegro-4}
\end{figure}
 \begin{figure}[H]
   \centering
    \includegraphics[width=0.8\linewidth]{graficos/capital_infinito/fibonacci_numeros_negros/corridas_5_f_i.png}
    \caption{Flujo de caja - 5 corridas}
    \label{fig:corridas_5_f_i_numnegro-4}
\end{figure}
 \begin{figure}[H]
   \centering
    \includegraphics[width=0.8\linewidth]{graficos/capital_infinito/fibonacci_numeros_negros/frecuencia_relativa_evolucion_f_i.png}
    \caption{Histograma de frecuencia}
    \label{fig:frecuencia_relativa_evolucion_f_i_numnegro-4}
\end{figure}
 \begin{figure}[H]
   \centering
    \includegraphics[width=0.8\linewidth]{graficos/capital_infinito/fibonacci_numeros_negros/frecuencias_multiples_f_i.png}
    \caption{Histograma de frecuencia}
    \label{fig:frecuencias_multiples_f_i_numnegro-4}
\end{figure}



\subsubsection{Martingala - Técnica - Números negros}
 \begin{figure}[H]
   \centering
    \includegraphics[width=0.8\linewidth]{graficos/capital_infinito/martingala_numeros_negros/flujo_caja_m_i.png}
    \caption{Flujo de caja - 1° corrida}
    \label{fig:flujo_caja_m_i_numnegro-4}
\end{figure}
 \begin{figure}[H]
   \centering
    \includegraphics[width=0.8\linewidth]{graficos/capital_infinito/martingala_numeros_negros/corridas_5_m_i.png}
    \caption{Flujo de caja - 5 corridas}
    \label{fig:corridas_5_m_i_numnegro-4}
\end{figure}
 \begin{figure}[H]
   \centering
    \includegraphics[width=0.8\linewidth]{graficos/capital_infinito/martingala_numeros_negros/frecuencia_relativa_evolucion_m_i.png}
    \caption{Histograma de frecuencia}
    \label{fig:frecuencia_relativa_evolucion_m_i_numnegro-4}
\end{figure}
 \begin{figure}[H]
   \centering
    \includegraphics[width=0.8\linewidth]{graficos/capital_infinito/martingala_numeros_negros/frecuencias_multiples_m_i.png}
    \caption{Histograma de frecuencia}
    \label{fig:frecuencias_multiples_m_i_numnegro-4}
\end{figure}



\subsubsection{Martingala Invertido - Técnica - Números negros}
 \begin{figure}[H]
   \centering
    \includegraphics[width=0.8\linewidth]{graficos/capital_infinito/martingala_invertida_numeros_negros/flujo_caja_o_i.png}
    \caption{Flujo de caja - 1° corrida}
    \label{fig:flujo_caja_mi_i_numnegro-4}
\end{figure}
 \begin{figure}[H]
   \centering
    \includegraphics[width=0.8\linewidth]{graficos/capital_infinito/martingala_invertida_numeros_negros/corridas_5_o_i.png}
    \caption{Flujo de caja - 5 corridas}
    \label{fig:corridas_5_mi_i_numnegro-4}
\end{figure}
 \begin{figure}[H]
   \centering
    \includegraphics[width=0.8\linewidth]{graficos/capital_infinito/martingala_invertida_numeros_negros/frecuencia_relativa_evolucion_o_i.png}
    \caption{Histograma de frecuencia}
    \label{fig:frecuencia_relativa_evolucion_mi_i_numnegro-4}
\end{figure}
 \begin{figure}[H]
   \centering
    \includegraphics[width=0.8\linewidth]{graficos/capital_infinito/martingala_invertida_numeros_negros/frecuencias_multiples_o_i.png}
    \caption{Histograma de frecuencia}
    \label{fig:frecuencias_multiples_mi_i_numnegro-4}
\end{figure}












\section{Conclusión}
A partir del análisis realizado, podemos afirmar que ninguna de las estrategias y técnicas clásicas de apuestas analizadas resulta favorable a largo plazo. Si bien pueden ofrecer ciertas ventajas momentáneas, especialmente en situaciones puntuales o bajo condiciones específicas, todas tienden a generar pérdidas sostenidas con el tiempo, ya sea con capital limitado o incluso con capital ilimitado.

Llama la atención que algunas estrategias parecieran ser apropiadas para una determinada técnica, pero muy arriesgadas para otra técnica. Como por ejemplo la estrategia "Martingala Invertida" para la técnica "Numero unico", en la cual hubo una sola bancarrota de las 5 corridas. Al utilizar "Martingala Invertida" para la técnica "Numeros negros" todas las corridas terminan en bancarrota.


No obstante, estas estrategias pueden ser consideradas según el perfil del jugador. Algunas resultan más apropiadas para apostadores conservadores, mientras que otras apelan a perfiles más arriesgados. Del mismo modo, jugadores con capitales elevados pueden sostener por más tiempo estrategias agresivas, aunque esto no garantice beneficios, sino solo una postergación de las pérdidas.

Las simulaciones y gráficos presentados permiten visualizar momentos clave donde conviene detenerse para asegurar alguna ganancia eventual o minimizar pérdidas.

Finalmente, aunque algunas estrategias implican un mayor riesgo con la promesa de mayores ganancias, en el contexto del casino ese riesgo es siempre elevado. La única certeza es que, en el largo plazo, la casa siempre tiene la ventaja. Identificar cuándo resistir y cuándo desistir se convierte así en la herramienta más útil del jugador racional.



\bibliographystyle{unsrt}  


\renewcommand{\refname}{Bibliografía}

\begin{thebibliography}{1}

\bibitem{referencia1}
Probabilidad y estadística con Python.\\
\url{https://relopezbriega.github.io/blog/2015/06/27/probabilidad-y-estadistica-con-python/}

\bibitem{referencia2}
Python para impacientes.\\
\url{https://python-para-impacientes.blogspot.com/2014/08/graficos-en-ipython.html}

\bibitem{referencia3}
Paul Meyer: Probabilidad y aplicaciones estadísticas. Addison - Wesley. Editorial Iberoamericana\\

\bibitem{referencia4}
Ruleta Online D'alembert.\\
\url{https://www.ruletaonline.org/estrategias/dalembert/}

\end{thebibliography}





\end{document}
